\scanpage{77}
\deleted{donc
\[
(\lambda^n-\lambda^{n'})g_{n'}=0\qquad\forall\lambda\in\mathbb Z^r
\]
et il suffit de voir que les $\lambda^n-\lambda^{n'}$ engendrent
l'idéal $1$, ou mieux}

si on exigeait qu'on a, pour tt $\mathbb Z$-algèbre $k$,
des compositions entre les $P_k(X,Y)\simeq P(X,Y)\otimes_{\mathbb Z}k$,
et en exigeant la décomposition plus haut
valable \deleted{pour}\inserted{dans} tt $P_k(\prod X_i,Y)$, avec $\lambda\in(k^\times)^{\unsure{r|n}}$. On
va mettre comme \uline{structure supplémentaire}
la \uuline{donnée} de la décomposition $(*)$, avec

e) Condition \unsure{des degrés|de degré} pour la composition
\diagram{\begin{tikzcd}[ampersand replacement=\&,column sep=small]
X_1 \& {\times\cdots\times} \& X_r\arrow[r,"{(n_1,\ldots,n_r)}"] \& Y\\
{Z^1_1\times\cdots Z^1_{s_1}}\arrow[u,"{m^1_1,\ldots,m^1_{s_1}}"'] \& \& {Z^r_1,\ldots,Z^r_{s_r}}\arrow[u,"{m^r_1,\ldots,m^r_{s_r}}"] \&
\end{tikzcd}}
% slanted
\margin{Oubli : les hom de la sous-catégorie $P_1$ sont bien ceux de degré $1$ pour la décomposition $P(X,Y)=\bigoplus_i P_i(X,Y)$}
% slanted; first symbols on two lines touch the scan boundary
\margin{faut aussi si $X_r=0$ $\unsure{P|\mathcal P}_{n_1,\ldots,n_r}(X_1,X_2,X_r\deleted{\unsure{,\ldots}};M)=\unsure{0}$ si $n_r\ne0$, $P_{n_1,\ldots,n_{r-1}}(X_1,\ldots,X_{r-1})$ si $n_r=0$}
le composé est de degré $(n_i m^i_1,n_i m^i_2,\ldots,n_i m^i_{s_i})$
en $Z^i_1,\ldots,Z^i_{s_i}$ ; de $=$

f) condition \inserted{\unsure{symétrie|de symétrie}} \deleted{\unsure{pour}} la permutation des facteurs

g) $\forall X\in\operatorname{Ob}P$, l'application canonique déduite de
$X\longrightarrow0$
\[
P_0(0,Y)\longrightarrow P_0(X,Y)
\]
est bijective (les ``applications constantes'' \deleted{\unsure{canoniques}}\inserted{de $X$ dans $Y$}
``ne dépendent pas de $X$'')

Si $k$ est un anneau comm. et si $P_1$ est muni
d'une structure $k$-linéaire par
\[
k\unsure{-|\longrightarrow}\operatorname{End}(P_1)\longleftarrow\operatorname{End}(P)
\]
\deleted{on dit que}
telle que $(*)$ plus haut soit vrai pour
$\lambda\in k^r$ (non seulement $\lambda\in\mathbb Z^r$), on dit que
$P$ est $k$-polynomiale.
% archive sheet number: 76
